\documentclass[11pt]{article}
\usepackage{explorer}
\newcommand{\ExplorerLanguage}{EN}
\newcommand{\ExplorerNumber}{06}
\newcommand{\ExplorerKind}{Worked solutions}
\begin{document}
\ExplorerTitle{High school · graphs and logic}{Routes without lifting the pencil}
\ExplorerSection{Exercise 1}
A graph has degrees 2,3,3,4. How many edges does it have?\par
The degree sum is 12, so there are 6 edges.\par
\ExplorerSection{Exercise 2}
Prove that the number of odd-degree vertices is even.\par
The degree sum is even. Even degrees do not affect parity, so the odd degrees must also sum to an even number. Thus there are an even number of them.\par
\ExplorerSection{Exercise 3}
Can a square and both diagonals be drawn without repeating an edge? Interior crossings are not vertices.\par
No: all four vertices have degree 3. Four odd-degree vertices prevent an Euler trail.\par
\ExplorerSection{Exercise 4}
Why do two disjoint cycles not admit one trail even though all degrees are even?\par
No edge connects the two components, so one cannot move between them. Connectivity is essential.\par
\ExplorerSection{Problem: complete reasoning}
A network has edges AB, BC, CD, DA, AC and CE, each 1 km long. Starting at A, visit every edge at least once and return to A. Which vertices have odd degree? What is the minimum route length? Give a route attaining it.\par
Degrees at A,B,C,D,E are 3,2,4,2,1, so A and E are odd. A closed route must add traversals joining these two vertices; their shortest connection is A--C--E, of length 2 km. The minimum is \(6 + 2 = 8\) km. Route A--B--C--D--A--C--E--C--A attains it: AC and CE are used twice and the other edges once.\par
\end{document}
